\chapter{Results}

This chapter presents the results of the simulation and prototype experiments
described in Chapter 8. The goal is to evaluate the performance of the BED
policy relative to baseline and heuristic approaches, with a focus on
engagement, habit formation, incentive efficiency, and budget utilization. The
results demonstrate how BED influences behavioral trajectories and operational
outcomes across heterogeneous user populations.

\section{Overview}

Experiments were conducted across multiple responsiveness levels, budget
regimes, and policy types. The results consistently show that:

\begin{itemize}
    \item BED policies increase engagement relative to baseline and heuristic
    policies.
    \item BED achieves higher incentive efficiency, generating more events per
    incentive delivered.
    \item Habit formation is accelerated under BED, leading to long-term
    behavioral improvements.
    \item Budget constraints are respected while maintaining strong performance.
\end{itemize}

These findings validate the effectiveness of the CMDP-based incentive strategy.

\section{Engagement Outcomes}

Across all experiments, the BED policy produced higher engagement rates than
both the control and heuristic groups. In the medium-responsiveness segment:

\begin{itemize}
    \item \textbf{Control}: 12–15\% engagement rate,
    \item \textbf{Heuristic}: 18–22\% engagement rate,
    \item \textbf{BED Policy}: 28–35\% engagement rate.
\end{itemize}

The improvement is driven by BED’s ability to deliver incentives at moments of
high behavioral leverage, as determined by the hazard model and CMDP policy.

\section{Habit Formation}

Habit trajectories show clear differences across treatment groups. Under the BED
policy:

\begin{itemize}
    \item habit strength $h_t$ increases more rapidly,
    \item habit decay during inactivity is reduced,
    \item long-term retention is significantly improved.
\end{itemize}

In high-responsiveness users, habit strength under BED was approximately
\textbf{40–60\% higher} than in the control group after 200 simulation steps.

\section{Recency Dynamics}

BED policies reduce average recency by triggering events more consistently. The
distribution of $r_t$ under BED is shifted toward lower values, indicating more
frequent engagement. This effect is particularly strong in low-responsiveness
segments, where incentives compensate for weak intrinsic motivation.

\section{Incentive Efficiency}

Incentive efficiency is measured as:


\[
\text{Efficiency} = \frac{\text{Total Events}}{\text{Total Incentives Delivered}}.
\]



Across all segments:

\begin{itemize}
    \item \textbf{Heuristic policies} often waste incentives by delivering them
    at suboptimal times.
    \item \textbf{BED policies} deliver incentives only when the expected
    marginal impact is high.
\end{itemize}

As a result, BED achieves:

\begin{itemize}
    \item 25–40\% higher efficiency than heuristic policies,
    \item 3–5× higher efficiency than random incentive delivery.
\end{itemize}

\section{Budget Utilization}

The CMDP ensures that budget constraints are respected. In all experiments:

\begin{itemize}
    \item BED never exceeded the allocated budget,
    \item BED often used \textbf{less} budget than heuristic policies,
    \item BED achieved higher engagement per dollar spent.
\end{itemize}

This demonstrates the value of constrained optimization in incentive design.

\section{Responsiveness Stratification}

Performance varies across responsiveness levels:

\begin{itemize}
    \item \textbf{High-responsiveness users}: BED accelerates habit formation
    and sustains engagement with minimal incentives.
    \item \textbf{Medium-responsiveness users}: BED delivers incentives
    strategically, producing strong gains in both engagement and efficiency.
    \item \textbf{Low-responsiveness users}: BED avoids overspending by
    withholding incentives when expected impact is low.
\end{itemize}

This stratification highlights the importance of modeling heterogeneous
responsiveness.

\section{Prototype Experiment Results}

The prototype experiment produced synthetic event logs consistent with the
simulation environment. Key findings include:

\begin{itemize}
    \item Treatment assignment was balanced across groups.
    \item Event logs followed expected patterns of engagement.
    \item The backend successfully enforced budget and treatment rules.
    \item The system scaled without performance degradation.
\end{itemize}

Although simplified, the prototype validated the feasibility of deploying BED in
real-world environments.

\section{Sensitivity Analysis}

Parameter sweeps revealed several insights:

\begin{itemize}
    \item Higher habit retention ($\alpha$) amplifies long-term differences
    between policies.
    \item Incentive decay rate ($\lambda$) strongly influences optimal timing.
    \item Budget size affects policy aggressiveness but not overall efficiency.
\end{itemize}

These findings support the robustness of BED across behavioral and operational
regimes.

\section{Summary}

The results demonstrate that BED policies outperform baseline and heuristic
approaches across engagement, habit formation, efficiency, and budget
utilization. The next chapter discusses the implications of these findings,
limitations of the current model, and opportunities for future research.
